\documentclass[11pt,letterpaper]{article}
\usepackage{fontspec}
\usepackage{tocloft}
\usepackage{multicol}
\usepackage{nicefrac}
\usepackage[left=1.3in,right=1.3in,top=1in,bottom=1in]{geometry}
\setmainfont[Scale=1.1,AutoFakeBold=1.5,AutoFakeSlant=0.3]{Doves Type}

% Prevent page breaks within list items
\widowpenalty=10000
\clubpenalty=10000
\interlinepenalty=500

\title{Pan-Seared Salmon with Lemon-Dill Butter Sauce}
\author{}
\date{}

\begin{document}

\maketitle
\thispagestyle{empty}

% Quote intentionally omitted

\section*{Ingredients}
\setlength{\columnsep}{20pt}
\begin{multicols}{2}
\noindent
    Salmon fillets, (6oz), skin-on \dotfill 4 \\
    Kosher salt \dotfill 1~tsp. \\
    Black pepper \dotfill \nicefrac{1}{2}~tsp. \\
    Unsalted butter \dotfill 6~Tbsp. \\
    Garlic cloves \dotfill 4 \\
    \columnbreak
    Dry white wine \dotfill \nicefrac{1}{2}~cup \\
    Lemon juice \dotfill \nicefrac{1}{4}~cup \\
    Lemon zest \dotfill 4~tsp. (from 2~lemons) \\
    Fresh dill \dotfill 3~Tbsp. \\
    Capers \dotfill 3~Tbsp. \\
    Fresh parsley \dotfill 3~Tbsp. \\
    Neutral oil \dotfill 1~Tbsp. \\
\end{multicols}

\section*{Directions}

\noindent
Thaw \textbf{salmon} if frozen ---
Mince \textbf{garlic}; set aside in \textit{Small Bowl~\#1} ---
Zest and juice \textbf{lemons}; combine \textbf{lemon zest}, \textbf{dill}, \textbf{capers}, and \textbf{parsley} in \textit{Medium Bowl~\#1} (sauce finish) ---
Measure \nicefrac{1}{4}~cup \textbf{lemon juice} and \nicefrac{1}{2}~cup \textbf{white wine}; have ready

\begin{enumerate}
    \item Pat \textbf{salmon} fillets dry with paper towels. Sprinkle both sides lightly with \textbf{kosher salt} (about 1~tsp. total for all 4) and \textbf{black pepper}. Rest at room temperature for \textit{15--20~minutes}. Pat dry again---surface moisture prevents crisp skin.

    \item Heat a nonstick skillet over medium-high until hot, about \textit{2~minutes}. Add a thin film of \textbf{neutral oil}. Working in 2 batches, place \textbf{salmon} fillets skin-side down (do not overcrowd) and press gently with a spatula for \textit{5~seconds}. Cook without moving for \textit{4--5~minutes} until skin is crisp and golden and the fillets release when nudged. Do not move the fillets during this time---moving prevents even browning. Continue cooking in \textit{30~second} increments if skin still sticks.

    \item Flip \textbf{salmon} and cook flesh-side down for \textit{2--3~minutes} until center reaches \textit{125--130°F} on an instant-read thermometer (medium, moist and flaky). Flesh should appear mostly opaque with a faint pink center; it will carry over \textit{5--10°F} during resting. Transfer to a platter and tent loosely with foil. Repeat with remaining fillets.

    \item Make the pan sauce in the same skillet (after all \textbf{salmon} is cooked):
\begin{enumerate}
    \item Reduce heat to medium. Add 3~Tbsp. \textbf{butter} and minced \textbf{garlic} (\textit{Small Bowl~\#1}). Cook, stirring, for \textit{30~seconds} until garlic is fragrant and golden---do not let it brown or burn.
    \item Add \nicefrac{1}{2}~cup \textbf{white wine} and simmer until reduced by half, about \textit{3--4~minutes}. The pan should have a thin, syrupy film of reduced wine.
    \item Add \nicefrac{1}{4}~cup \textbf{lemon juice}, then \textbf{lemon zest}, \textbf{dill}, \textbf{capers}, and \textbf{parsley} (\textit{Medium Bowl~\#1}). Remove from heat and swirl in remaining 3~Tbsp. \textbf{butter} until melted and glossy. Taste and adjust \textbf{salt} if needed. Spoon sauce generously over \textbf{salmon} and serve immediately.
\end{enumerate}
\end{enumerate}

\newpage

{\footnotesize
\setlength{\columnsep}{20pt}
\setlength{\multicolsep}{6pt}
\begin{multicols}{2}
\setlength{\parindent}{0pt}
\setlength{\parskip}{2pt}
\setlength{\itemsep}{0pt}
\setlength{\parsep}{0pt}

\subsection*{Yield}
\begin{itemize}
    \item Serves 4 (4 fillets, 6~oz. each)
\end{itemize}

\subsection*{Equipment Required}
\begin{itemize}
    \item 12-inch nonstick skillet (or cook in 2 batches in a 10-inch)
    \item Instant-read thermometer
    \item Spatula (thin, flexible preferred for flipping)
    \item Platter for holding cooked \textbf{salmon}
    \item Small prep bowl (1), Medium prep bowl (1)
    \item Paper towels
    \item Measuring cups and spoons
\end{itemize}

\subsection*{Mise en Place}
\begin{itemize}
    \item Small Bowl \#1 --- minced \textbf{garlic}
    \item Medium Bowl \#1 --- sauce finish: \textbf{lemon zest}, \textbf{dill}, \textbf{capers}, \textbf{parsley}
    \item Measure \nicefrac{1}{4}~cup \textbf{lemon juice} and \nicefrac{1}{2}~cup \textbf{white wine}; have ready
    \item Thaw \textbf{salmon} fully before cooking; allow \textit{15--20~minutes} at room temperature after dry brine
\end{itemize}

\subsection*{Ingredient Tips}
\begin{itemize}
    \item \textbf{Salmon}: Wild or farmed both work; skin-on is essential for crisp texture
    \item \textbf{White wine}: Dry varieties (Pinot Grigio, Sauvignon Blanc) work best
    \item \textbf{Capers}: Rinse if you prefer less salt; chop if large
    \item \textbf{Frozen salmon}: Thaw overnight in fridge or in sealed bag in cold water \textit{30--45~minutes}
\end{itemize}

\subsection*{Preparation Tips}
\begin{itemize}
    \item Dry brine draws moisture from the surface---pat dry again before searing so skin crisps instead of steaming
    \item Do not move \textbf{salmon} during skin-side cooking; patience yields crisp, release-ready skin
    \item Pull \textbf{salmon} at \textit{125--130°F} for moist, flaky texture; carryover will raise it \textit{5--10°F}. USDA recommends \textit{145°F} for safety; commercially sourced salmon is generally considered safe at lower temps
    \item Pan sauce comes together quickly---have \textbf{garlic} and \textit{Medium Bowl~\#1} ready before you start the sear
    \item If skin sticks, it is not ready; give it another \textit{30--60~seconds} before flipping
\end{itemize}

\subsection*{Make Ahead \& Storage}
\begin{itemize}
    \item Dry brine can start \textit{15--20~minutes} before cooking
    \item Sauce ingredients can be prepped ahead; make sauce immediately before serving
    \item Leftover \textbf{salmon} keeps \textit{2--3~days} refrigerated; reheat gently to avoid drying
    \item Sauce does not reheat well---make fresh
\end{itemize}

\subsection*{Serving Suggestions}
\begin{itemize}
    \item Serve immediately while skin is crisp and sauce is warm
    \item Pairs well with roasted asparagus, cheesy scalloped potatoes, or rice
    \item A squeeze of fresh \textbf{lemon} at the table adds brightness
\end{itemize}

\end{multicols}
}

\end{document}
